\chapter*{Conventions typographiques}
\markboth{CONVENTIONS TYPOGRAPHIQUES}{}
\addcontentsline{toc}{chapter}{Conventions typographiques}

\section*{Dans le texte}

\begin{description}[leftmargin=3.2cm, style=nextline]
  \item[\texttt{fonte \`a chasse fixe}] Nom de fichier, commande \`a saisir,
        variable d'environnement, adresse d'une page.
  \item[\textbf{Gras}] \'El\'ement de l'interface sur lequel agir : nom
        d'onglet, libell\'e de bouton, intitul\'e de champ.
  \item[\textit{Italique}] Terme d\'efini au glossaire, ou insistance.
  \item[EF-nn, CR-nn] Renvoi \`a une exigence fonctionnelle ou \`a un
        crit\`ere de recette du cahier des charges. Ils sont rappel\'es ici
        pour que le lecteur puisse remonter au document d'origine.
\end{description}

\section*{Encadr\'es}

Trois encadr\'es rythment le guide. Ils ne sont pas d\'ecoratifs : chacun
signale une nature d'information diff\'erente.

\begin{tcolorbox}[colback=green!4!white, colframe=darkgreen,
  title=\textbf{Bon \`a savoir}, fonttitle=\bfseries, boxrule=0.6pt]
Pr\'ecision qui \'eclaire un comportement de la plateforme, sans que son
ignorance ne porte \`a cons\'equence.
\end{tcolorbox}

\begin{tcolorbox}[colback=orange!5!white, colframe=orange!70!black,
  title=\textbf{Attention}, fonttitle=\bfseries, boxrule=0.6pt]
Point sur lequel une erreur de manipulation a un effet visible : donn\'ees
perdues, service interrompu, r\'esultat trompeur.
\end{tcolorbox}

\begin{tcolorbox}[colback=red!4!white, colframe=darkred,
  title=\textbf{Limitation connue}, fonttitle=\bfseries, boxrule=0.6pt]
Comportement que la plateforme n'a pas, ou qu'elle a de fa\c{c}on incompl\`ete,
et qui est \'enonc\'e ici plut\^ot que d\'ecouvert en exploitation.
\end{tcolorbox}

\section*{Commandes}

Les commandes \`a saisir sont pr\'ec\'ed\'ees d'une invite qui indique le
contexte :

\begin{lstlisting}[language=bash]
$ commande lancee par un utilisateur ordinaire
# commande necessitant les droits d'administration
PS> commande PowerShell, sous Windows
\end{lstlisting}

L'invite elle-m\^eme ne se saisit pas.
